\documentclass[a4paper, oneside]{article}
\usepackage{graphicx}
\usepackage{amsthm}
\usepackage{amsmath}
\usepackage{amssymb}
\usepackage[a4paper,
            bindingoffset=0.2in,
            left=2cm,
            right=2cm,
            top=2cm,
            bottom=2cm,
            footskip=.25in]{geometry}
\usepackage[italian]{babel}
\usepackage{pgfplots}
\usepackage{tabularx}
\usepackage{tikz}
\usepackage{wrapfig}
\usepackage{color}
\usepackage[d]{esvect}
\definecolor{page}{rgb}{0.129,0.157,0.212}
\pagecolor{page}
\color{white}
\graphicspath{ {./images/} }
\usetikzlibrary{shapes.geometric}
\usetikzlibrary{datavisualization}
\usetikzlibrary{datavisualization.formats.functions}
\usetikzlibrary{patterns}
\pgfplotsset{width=10cm,compat=1.18}

\title{Appunti di Metodi}
\author{Tommaso Miliani}
\date{13-03-26}

\begin{document}
\newtheoremstyle{theoremEnv}
                {}          % Space above
                {}          % Space below
                {\slshape}  % Body font
                {}          % Indent amount
                {\bfseries} % Head font
                {.}         % Punctuation after head
                {\newline}  % Space after theorem head
                {}          % Theorem head spec
\theoremstyle{theoremEnv}

\newtheorem{definition}{Definizione}[section]
\newtheorem{theorem}{Teorema}[section]
\newtheorem{lemma}{Proposizione}[section]
\newtheorem{observation}{Osservazione}[section]
\newtheorem{corollary}{Corollario}[theorem]
\newtheorem{example}{Esempio}[section]
\newtheorem{remark}{Enunciato}[section]

\maketitle

\section{Formula integrale di Cauchy}
La formula integrale di Cauchy
\begin{align}
    f(z) = \frac{1}{2\pi i} \int_{\Gamma}^{} \frac{f(w)}{w - z} \ dw 
\end{align}
Dove la funzione deve essere Olomorfa in tutti i punti del circuito. Le conseguenze
di questa formula sono le seguenti

\begin{corollary}[Conseguenza della formula integrale]
    Una funzione olomorfa è necessariamente $C^{\infty }$, inoltre 
    le derivata possono essere calcolate con una formula simile:
    \begin{align}
        f^{(n)}(z) = \frac{n!}{2\pi i} \int_{\Gamma} \frac{f(w)}{(w - z)^{ n - 1}}\ dw 
    \end{align}
    Con $\Gamma$ lo stesso circuito di quello scelto per la funzione base. La dimostrazione si ottiene derivando 
    la formula integrale di Cauchy.
\end{corollary}

\begin{theorem}[Teorema di analicità di una funzione]
    Sia $f$ una funzione olomorfa definita su un aperto $\mathbb{A}$ , sia $z_0  \in \mathbb{A}$, 
    tale che 
    \begin{gather*}
        B=(z_0, r ) \subset \mathbb{A} \qquad r > 0
    \end{gather*}
    Allora la funzione può essere espansa in serie
    \begin{align}
        f(z) = f(z_0) + f'(z_0) (z - z_0) + \frac{f^{(n)}(z_0)}{n!}(z - z_0)^{n}
    \end{align}
    E converge sul disco $B$. Esiste sempre un intorno non banale in cui 
    la serie di Taylor converge.
\end{theorem}
\begin{proof}
    (Non la chiede). Si utilizza uno sviluppo in serie all'interno della 
    forma integrale di Cauchy, si sviluppa a partire dalla palla $B$, e si prende
    la frontiera $\Gamma$ (che sta fuori dal disco in quanto è aperto), e si sceglie $z$ 
    dentro la frontiera:
    \begin{gather*}
        \frac{1}{w - z} = \frac{1}{(w - z_0) - (z - z_0)} = \frac{1}{w - z_0} \frac{1}{1 - \frac{z - z_0}{w - z_0}}
    \end{gather*}
    Utilizzando la disuguaglianza triangolare
    \begin{gather*}
        \left| z - z_0 \right| < \left| w - z_0 \right|  
    \end{gather*}
    Se si vede 
    \begin{gather*}
        \frac{1}{1 - \frac{z - z_0}{w - z_0}}
    \end{gather*}
    Come serie geometrica convergente, allora essa è convergente 
    perché
    \begin{gather*}
        \left| \frac{z - z_0}{w - z_0} \right| < 1
    \end{gather*}
    E dunque 
    \begin{gather*}
        \frac{1}{w - z_0} \sum_{n = 0}^{\infty } \left(\frac{z - z_0}{w - z_0}\right) ^{n}
    \end{gather*}
    Allora si può scrivere
    \begin{gather*}
        f(z) = \frac{1}{2\pi i} \int_{\Gamma} \frac{f(z)}{w - z} dw = \frac{1}{2\pi i} \int_{\Gamma} \frac{f(w)}{w - z_0} \sum_{n = 0}^{\infty } \left(\frac{z - z_0}{w - z_0}\right)^{n}dq
    \end{gather*}
    Si può allora sostituire l'integrale con la serie
    \begin{gather*}
        = \sum_{n = 0}^{\infty } \frac{(z - z_0)^{n}}{2\pi i} \int_{\Gamma} \frac{f(w)}{(w - z_0)^{n + 1}} \ dw  = \sum_{n = 0}^{\infty } \frac{f^{(n)}(z_0)}{n!}(z - z_0)^{n} 
    \end{gather*}
    Se si considera
    \begin{gather*}
        \sup_{\Gamma} \left| \frac{f(w)}{w - z_0} \left(\underset{\left| \delta \right|  < 1}{\underline{\frac{z - z_0}{w - z_0}}}\right)^{n} \right|  \leq \dfrac{\sup_{\Gamma} \left| f(w) \right| }{r}\left| \delta \right| ^n
    \end{gather*}
    Il termine sulla frazione in fondo  è un termine di una serie convergente e duneque la serie 
    converge totalmente (in quanto $\delta$ converge).
\end{proof}

\begin{theorem}[Condizioni di olomorfismo]
    Se $f \in \mathbb{C}^{0}(\mathbb{A}, \mathbb{C})$, ossia una funzione continua su 
    di un aperto $\mathbb{A} \subset \mathbb{C}$. Se essa è continua e ammette primitiva, 
    allora è olomorfa su $\mathbb{A}$.
\end{theorem}
\begin{proof}
    Se $f = F'$, dunque se $F$ è olomorfa, allora anche la derivata di 
    una funzione olomorfa è olomorfa (si vede sviluppando in serie), 
    allora $f$ è olomorfa.
\end{proof}

\begin{theorem}[Riassunto delle proprietà di funzioni continue su di un aperto]
    Sia $f$ una funzione definita su di un aperto $\mathbb{A}$ 
    $f : \mathbb{A} \to \mathbb{C}$, se essa è continua, allora 
    \begin{enumerate}
        \item $f$ è olomorfa
        \item $f$ è localmente esatta: così come nel caso reale,
        l'integrale su qualsiasi circuito in $\mathbb{A}$ è nullo.
        \item $f$ ammette localmente primitiva
        \item 
        \begin{align}
            \oint_\gamma f \ dz = 0
        \end{align}
        per ogni circuito $\gamma$ nullomotopo in $\mathbb{A}$.
    \end{enumerate}
\end{theorem}
\begin{corollary}
    Il dominio della funzione deve essere semplicemente connesso per far sì che valga 
    la condizione 3, per questo vale \textbf{localmente} e non globalmente. Un esempio 
    è la funzione
    \begin{gather*}
        f(z) = \frac{1}{z - z_0} \ \Longrightarrow \ F = \log(z - z_0)
    \end{gather*}
    Essa è una buona primitiva localmente ma non globalmente, infatti il dominio presenta un taglio
    in $z = z_0$. 
\end{corollary}


\section{Serie di Laurent}
\begin{definition}
    Si definisce \textbf{serie bilatera}  una serie
    \begin{align}
        \sum_{n = -\infty }^{\infty } c_n(z - z_0)^{n} 
    \end{align}
    ossia che una serie che ha anche potenze negative. Si è interessati 
    a vedere come converge questa tipologia di serie. La convergenza
    di questa tipologia di serie viene studiata mediante due convergenze:
    \begin{align}
        \left\{\begin{array}{l}
            \text{parte regolare} \qquad \sum_{n = 0}^{\infty } c_n(z - z_0)^{n} \\
            \text{singolare} \qquad \sum_{n = -\infty }^{-1} c_n(z - z_0)^{n} 
        \end{array}\right.
    \end{align}
    Dove la parte regolare è la semplice serie di Taylor, che converge in una
    palla definita da un raggio $r_2$, per cui si associa un raggio di convergenza
    \begin{align}
        \left| z - z_0 \right| < r_2
    \end{align}
    Per far convergere la serie singolare si può considerare:
    \begin{align}
        w = \frac{1}{z - z_0} \ \Longrightarrow \ \sum_{n = 1}^{\infty } c_{-n}w^{n}
    \end{align}
    Dove la convergenza si ha considerando:
    \begin{align}
        \left| w \right| < \rho \Longleftrightarrow  \left| z - z_0 \right|   > \frac{1}{\rho} \equiv r_1 \qquad \qquad w = \frac{1}{\left| z - z_0 \right| }
    \end{align}
    Dunque la serie singolare converge all'esterno del disco, mentre la parte regolare converge all'esterno di un disco. Esistono dunque 
    tre casi di convergenza
    \begin{enumerate}
        \item $r_1 > r_2$: non esiste nessun dominio in cui vi è convergenza di entrambi, allora 
        la serie bilatera non converge.
        \item $r_1 = r_2$: la serie può al più convergere sul bordo 
        del disco (ma non si hanno condizioni forti per la convergenza), su qualche 
        punto può convergere per $\left| z - z_0  \right| = r_1 = r_2$.
        \item $r_1 < r_2$: la convergenza si ha su di una corona circolare 
        \begin{align}
            \mathbb{B} = (z_0, (z_1, z_2))
        \end{align}
    \end{enumerate}
\end{definition} \noindent
È possibile dunque utilizzare questa tipologia di serie per ottenere
funzioni olomorfe.
\begin{theorem}[Sviluppo in serie di Laurent]
    Presa una funzione $f$ olomorfa in un dominio $\mathbb{A}$ aperto, e si assume che
    esista una corona circolare
    \begin{gather*}
        \mathbb{B} = (z_0, (r_1, r_2)) \subset \mathbb{A}
    \end{gather*}
    Ovviamente $0 \leq r_1 < r_2 \leq + \infty$, allora 
    \begin{align}
        \forall z \in \mathbb{B}, \qquad f(z) = \sum_{n = -\infty }^{\infty } c_n(z - z_0)^{n} 
    \end{align}
    e converge assolutamente e totalmente. Inoltre si può scrivere
    \begin{align}
        c_n = \frac{1}{2\pi i} \int_{\gamma_r} \frac{f(w)}{(w - z_0)^{ n + 1}} \ dw  
    \end{align}
    Dove $\gamma_r$ è una qualunque circonferenza contenuta completamente nella corona $\mathbb{B}$, 
    tale che $\gamma_r : z_0 + re^{i\theta} \ r_1 < r < r_2$.
\end{theorem}

\begin{corollary}
    Questo teorema consente di far convergere la serie anche intorno ad un buco 
    (l'importante è la corona sia totalmente contenuta nel dominio di olomorfia).
\end{corollary}

\section{Singolarità isolate}
Sia $f$ una funzione definita su di un aperto $\mathbb{A}$, posto 
$\mathcal{N} = \{a_i\}$ un insieme di punti singolari in cui la funzione non è 
definita all'interno di $\mathbb{A}$. La funzione 
$f \in H(\mathbb{A} \ \mathcal{N})$, le singolarità prendono il nome
di \textbf{isolate} se non esiste un punto di accumulazione per le singolarità 
(ossia tale per cui per ogni intorno esista alemeno un altro punto dell'insieme
che ci cade dentro). Si definiscono diversi tipi di singolarità isolate:
\begin{itemize}
    \item  \textbf{singolarità eliminabile}: è una singolarità
isolata tale per cui la funzione può essere estesa anche a questa singolarità. La serie bilatera
ha solo potenze non negative. Se, sviluppando in serie di Laurent si ottiene che
la singolarità non esiste, allora essa è eliminabile e la funzione è estendibile
sulla singolarità $a$, definendo $f(a_i) = c_0$.
    \item \textbf{polo}: una singolarità isolata in cui la serie singolare ha 
    un numero finito di termini. 
    \begin{align}
        \exists m > 0 : \ c_{-m} \neq 0, \quad c_{-n} = 0 \ \forall n > m
    \end{align}
    Si dice allora che il polo è di \textbf{ordine} $m$: in un intorno del polo
    la funzione va come 
    \begin{gather*}
        f \approx \frac{1}{(z - a)^{m}}
    \end{gather*}
    \item \textbf{singolarità essenziale}: essa corrisponde ad al caso in 
    cui la parte singolare ha un numero infinito di termini. 
\end{itemize}
\begin{example}[Singolarità]
    \begin{itemize}
        \item Eliminabile: la funzione $\frac{\sin(z)}{z}$, 
        essa è estendibile 
        \item Polo: $\frac{1}{(z - a)^{m}}$, ossia un polo di ordine $m$.
        \item Essenziale: $\exp\left(-\frac{1}{z^{2}}\right)$ intorno a zero 
        può essere espanso in serie di Laurent che convergerà per $z \neq 0$:
        \begin{gather*}
            \sum_{n = 0}^{\infty } \frac{1}{n!} \left(- \frac{1}{z^{2}}\right) ^{m} = \sum_{n = -\infty }^{0} \frac{1}{\left| n \right|! } \left(-z^{2}\right)^{m}
        \end{gather*}
        L'origine è dunque un punto singolare.
    \end{itemize}
\end{example}

\begin{theorem}
    Considerata $a$ una singolarità isolata di una funzione $f$ olomorfa, sono equivalenti
    \begin{enumerate}
        \item $a$ è eliminabile
        \item $\lim_{z \to a} f(z)$ esiste su $\mathbb{C}$:
        \item $f$ è limitata in un intorno di $a$ (con $a$ escluso). 
    \end{enumerate}
    Invece in prossimità di un polo (ma non di una singolarità essenziale) il comportamento di
$f$ è divergente, cioè il modulo di $f$ cresce in modo illimitato.
\end{theorem}

\begin{theorem}
    $a$ è un polo se e solo se
    \begin{align}
        \lim_{z \to a} f(z) = \infty \Longleftrightarrow \lim_{z \to a} \left| f(z) \right| = +\infty 
    \end{align}
\end{theorem}

\begin{theorem}[Teorema di Casorati-Weirestrass]
    Se $a$ è singolarità essenziale se e solo se 
    \begin{align}
        \forall \delta > 0 : \ \mathbb{B} = (a, \delta) \subset \mathbb{A} \qquad f(\mathbb{B} \ {a})
    \end{align}
    dove $\mathbb{A}$ è un aperto, l'immagine di $f$ è \textbf{densa} in $\mathbb{C}$.
\end{theorem}

\begin{theorem}{Teorema di Picard}
    In ogni introno di una singolarità essenziale  $f$ assume tutti i valori
    di $\mathbb{C}$, tranne che al più un valore, infinite volte.   
\end{theorem}

\begin{example}
    La funzione 
    \begin{gather*}
        f(z) = e^{-\frac{1}{z}}
    \end{gather*}
    non assume mai valore in zero. 
\end{example}

\subsection{Trovare gli sviluppi di Laurent negli intorni delle singolarità}
Si possono definire sviluppi di Laurent in diversi insiemi. Se si prende una qualunque
funzione olomorfa, è posibile considerare diverse palle per poter determinare 
l'espansione di Laurent, così che lo sviluppo in serie non sia lo stesso:
lo sviluppo dipende dunque a seconda del punto di origine e dalla palla scelta.
\begin{example}[Esempio di funzione espandibile in modi diversi]
    Si considera lo sviluppo di 
    \begin{gather*}
        f(z) = \frac{1}{z(1 - z)}
    \end{gather*}
    Essa ha due singolarità $z = 0, z = 1$ sull'asse reale. Se ci si chiedesse 
    quali sviluppi di Laurent esistano fuori dall'origine:
    \begin{gather*}
        \mathbb{B}_1 = (0, (0, 1)) \qquad \mathbb{B}_2(0, (1, \infty ))
    \end{gather*}
    Nel primo caso 
    \begin{gather*}
        \frac{1}{z}\frac{1}{z - 1} = \frac{1}{z} \sum_{n = 0}^{\infty } z^{n} = \sum_{m = 1}^{\infty } z^{m}  
    \end{gather*}
    L'origine è una singolarità che è un polo di ordine 1. Nel secondo caso invece
    \begin{gather*}
        \frac{1}{z} \frac{1}{z - 1} = \frac{1}{z} \frac{1}{z\left(\frac{1}{z} - 1\right)} = -\frac{1}{z^{2}} \frac{1}{1 - \frac{1}{z}}
    \end{gather*}
    Si espande dunque l'ultima frazione 
    \begin{gather*}
        - \frac{1}{z^{2} \sum_{n = 0}^{\infty } } \left(\frac{1}{z}\right)^{n} = -\sum_{n = 0}^{\infty } \left(\frac{1}{z}\right) ^{n + 2}
    \end{gather*}
    Questi due sviluppi sono ovviamente diversi
\end{example}

\end{document}